% =====================================================================
%  Tablas, listas y cajas — suite de documentos LaTeX
%  Nivel 2 (Capacidades) · clase article.
%  Compilar: latexmk -pdf tablas-listas-cajas.tex
% =====================================================================
\documentclass[11pt,a4paper]{article}

% --- Base estándar ---
\usepackage[utf8]{inputenc}
\usepackage[T1]{fontenc}
\usepackage{lmodern}
\usepackage{microtype}
\usepackage[spanish,es-noshorthands]{babel}
\usepackage{csquotes}
\usepackage{amsmath}\usepackage{amssymb}
\usepackage[margin=2.4cm]{geometry}
\usepackage{xcolor}
\usepackage{listings}
% --- Tablas ---
\usepackage{booktabs}
\usepackage{tabularx}
\usepackage{multirow}
\usepackage{siunitx}
\usepackage{longtable}
% --- Listas y cajas ---
\usepackage{enumitem}
\usepackage[most]{tcolorbox}
% --- Enlaces ---
\usepackage{hyperref}
\hypersetup{unicode, pdftitle={Tablas, listas y cajas en LaTeX},
            pdfauthor={Rodrigo Martín Vidal Galán},
            colorlinks=true, linkcolor=blue!50!black, urlcolor=blue!50!black}
\usepackage[spanish,capitalise,nameinlink]{cleveref}

% --- Estilo de código ---
\definecolor{codeback}{RGB}{245,247,250}
\definecolor{codekw}{RGB}{20,60,120}
\definecolor{codecom}{RGB}{110,120,135}
\lstset{
  basicstyle=\ttfamily\footnotesize,
  keywordstyle=\color{codekw}\bfseries,
  commentstyle=\color{codecom}\itshape,
  backgroundcolor=\color{codeback},
  frame=single, rulecolor=\color{codekw!30},
  breaklines=true, showstringspaces=false, columns=fullflexible,
  language=[LaTeX]TeX, xleftmargin=8pt, framesep=4pt, extendedchars=true,
  literate={á}{{\'a}}1 {é}{{\'e}}1 {í}{{\'i}}1 {ó}{{\'o}}1 {ú}{{\'u}}1
           {ñ}{{\~n}}1 {¿}{{?`}}1 {¡}{{!`}}1,
}
\newcommand{\C}[1]{\texttt{\textbackslash #1}}
\newcommand{\pkg}[1]{\textsf{#1}}
\newcommand{\produce}{\par\smallskip\noindent\textit{produce:}\par\smallskip}

\title{\textbf{Tablas, listas y cajas en \LaTeX}\\[2pt]
       \large Elementos de contenido}
\author{Rodrigo Martín Vidal Galán}
\date{Junio de 2026}

\begin{document}
\maketitle

\noindent Los tres elementos que estructuran el contenido \emph{dentro} del
texto: \textbf{tablas} (datos en filas y columnas), \textbf{listas} (ítems
ordenados o no) y \textbf{cajas} (encuadrar y resaltar). Formato código →
resultado. \emph{(El canon tipográfico está en \emph{Estándares}; este documento
lo aplica y amplía.)}

\newpage
\tableofcontents
\newpage

% =====================================================================
\section{Tablas: lo básico}
% =====================================================================
Una tabla se arma con el entorno \texttt{tabular} y un \textbf{preámbulo de
columnas} que dice cuántas hay y cómo se alinean:
\begin{itemize}[noitemsep,leftmargin=*]
  \item \texttt{l} \texttt{c} \texttt{r}: columna alineada a izquierda, centro o
        derecha (ancho automático).
  \item \texttt{p\{3cm\}}: columna de ancho fijo con texto que \textbf{ajusta}
        (varias líneas).
  \item \texttt{\&} separa celdas; \texttt{\textbackslash\textbackslash} termina
        la fila.
\end{itemize}
\begin{lstlisting}
\begin{tabular}{lcr}
  Nombre & Edad & Nota \\
  Ana    & 21   & 9.5  \\
  Beto   & 23   & 8.0  \\
\end{tabular}
\end{lstlisting}

% =====================================================================
\section{El canon: \texttt{booktabs}}
% =====================================================================
Las tablas profesionales usan \pkg{booktabs}: solo reglas \textbf{horizontales}
(\C{toprule}, \C{midrule}, \C{bottomrule}), \textbf{nunca verticales}. El aire y
las líneas finas se leen mucho mejor que una grilla cerrada.
\begin{lstlisting}
\begin{tabular}{@{}lcr@{}}
  \toprule
  Nombre & Edad & Nota \\
  \midrule
  Ana  & 21 & 9.5 \\
  Beto & 23 & 8.0 \\
  \bottomrule
\end{tabular}
\end{lstlisting}
\produce
\begin{center}
\begin{tabular}{@{}lcr@{}}
  \toprule
  Nombre & Edad & Nota \\
  \midrule
  Ana  & 21 & 9.5 \\
  Beto & 23 & 8.0 \\
  \bottomrule
\end{tabular}
\end{center}
\noindent El \texttt{@\{\}} al principio y final quita el relleno lateral
sobrante. Para una regla \textbf{parcial} (bajo algunas columnas) se usa
\C{cmidrule\{2-3\}}.

% =====================================================================
\section{Combinar celdas: \texttt{multicolumn} y \texttt{multirow}}
% =====================================================================
\begin{lstlisting}
\begin{tabular}{@{}lcc@{}}
  \toprule
  & \multicolumn{2}{c}{Parciales} \\   % una celda que ocupa 2 columnas
  \cmidrule(lr){2-3}
  Alumno & 1.\textordmasculine{} & 2.\textordmasculine{} \\
  \midrule
  \multirow{2}{*}{Grupo A} & 7 & 8 \\  % una celda que ocupa 2 filas
                           & 9 & 6 \\
  \bottomrule
\end{tabular}
\end{lstlisting}
\produce
\begin{center}
\begin{tabular}{@{}lcc@{}}
  \toprule
  & \multicolumn{2}{c}{Parciales} \\
  \cmidrule(lr){2-3}
  Alumno & 1.\textordmasculine{} & 2.\textordmasculine{} \\
  \midrule
  \multirow{2}{*}{Grupo A} & 7 & 8 \\
                           & 9 & 6 \\
  \bottomrule
\end{tabular}
\end{center}

% =====================================================================
\section{Anchos y alineación de números}
% =====================================================================
\paragraph{\texttt{tabularx}: ancho total fijo.} La columna \texttt{X} se
\textbf{estira} para que la tabla ocupe el ancho pedido (ideal para texto largo):
\begin{lstlisting}
\begin{tabularx}{\linewidth}{@{}lX@{}}
  \toprule
  Campo & Descripcion \\
  \midrule
  Titulo & Un texto que puede ser largo y se acomoda solo en varias lineas. \\
  \bottomrule
\end{tabularx}
\end{lstlisting}
\produce
\begin{center}
\begin{tabularx}{\linewidth}{@{}lX@{}}
  \toprule
  Campo & Descripción \\
  \midrule
  Título & Un texto que puede ser largo y se acomoda solo en varias líneas sin
           que tengas que cortarlo a mano. \\
  \bottomrule
\end{tabularx}
\end{center}

\paragraph{Columna \texttt{S} de \texttt{siunitx}: alinear por el decimal.}
\begin{lstlisting}
\begin{tabular}{@{}lS@{}}
  \toprule Magnitud & {Valor} \\ \midrule
  Ancho  & 12.5  \\
  Alto   &  3.25 \\
  Largo  & 120.0 \\
  \bottomrule
\end{tabular}
\end{lstlisting}
\produce
\begin{center}
\begin{tabular}{@{}lS@{}}
  \toprule Magnitud & {Valor} \\ \midrule
  Ancho  & 12.5  \\
  Alto   &  3.25 \\
  Largo  & 120.0 \\
  \bottomrule
\end{tabular}
\end{center}
Los números quedan alineados \textbf{por la coma decimal} (el encabezado va entre
\texttt{\{\}} para que no lo interprete como número).

% =====================================================================
\section{Tablas flotantes y muy largas}
% =====================================================================
Para que una tabla tenga \textbf{número y epígrafe} y se referencie, va en un
entorno \texttt{table} (flotante, como las figuras):
\begin{lstlisting}
\begin{table}[htbp]\centering
  \begin{tabular}{...} ... \end{tabular}
  \caption{Resultados.}\label{tab:res}   % caption ANTES del label
\end{table}
... y se cita con \cref{tab:res}.
\end{lstlisting}
Si la tabla \textbf{no entra en una página}, \texttt{table} no sirve (no se
parte): usá \pkg{longtable}, que se corta entre páginas repitiendo el encabezado.

% =====================================================================
\section{Listas}
% =====================================================================
Tres entornos básicos:
\begin{lstlisting}
\begin{itemize} \item ... \end{itemize}      % vinetas
\begin{enumerate} \item ... \end{enumerate}  % numeradas
\begin{description} \item[Termino] ... \end{description}  % con etiqueta
\end{lstlisting}
\produce
\begin{center}\begin{minipage}{0.85\linewidth}
\begin{description}[leftmargin=2.2cm,style=nextline,font=\bfseries]
  \item[itemize] lista con viñetas, para ítems sin orden.
  \item[enumerate] numerada, cuando el orden importa.
  \item[description] cada ítem tiene una \textbf{etiqueta} en negrita.
\end{description}
\end{minipage}\end{center}

\paragraph{Personalizar con \texttt{enumitem}.} Permite ajustar espaciado,
etiquetas y reanudar la numeración, sin \enquote{hacks}:
\begin{lstlisting}
\begin{enumerate}[label=(\alph*), noitemsep]  % (a) (b) (c), compacta
\begin{itemize}[label=$\Rightarrow$, leftmargin=*]
\begin{enumerate}[resume]   % CONTINUA la numeracion de una lista anterior
\end{lstlisting}
\produce
\begin{center}\begin{minipage}{0.7\linewidth}
\begin{enumerate}[label=(\alph*), noitemsep]
  \item primer ítem, etiqueta \texttt{(a)};
  \item segundo, sin espacio extra entre ítems;
  \item tercero.
\end{enumerate}
\end{minipage}\end{center}
Las listas \textbf{se anidan} hasta cuatro niveles (cambia la viñeta sola). Para
una lista \textbf{en línea} (dentro de un párrafo) está el entorno \texttt{enumerate*}
(de \texttt{enumitem} con la opción \texttt{inline}).

% =====================================================================
\section{Cajas}
% =====================================================================
\paragraph{Cajas simples (núcleo de \LaTeX).}
\begin{lstlisting}
\fbox{texto}                 % recuadro ajustado
\framebox[5cm][c]{centrado}  % recuadro de ancho fijo
\colorbox{yellow}{resaltado} % fondo de color
\parbox{4cm}{varias lineas}  % caja de ancho fijo con parrafo
\end{lstlisting}
\produce
\begin{center}
\fbox{texto} \quad \colorbox{yellow}{resaltado} \quad
\framebox[3cm][c]{ancho fijo}
\end{center}

\paragraph{\texttt{minipage}: una \enquote{mini página}.} Encierra contenido en un
ancho dado; es la base para poner cosas \textbf{lado a lado}:
\begin{lstlisting}
\begin{minipage}{0.48\textwidth} ... izquierda ... \end{minipage}\hfill
\begin{minipage}{0.48\textwidth} ... derecha ...   \end{minipage}
\end{lstlisting}

\paragraph{Cajas vistosas con \texttt{tcolorbox}.} Para destacados, definiciones o
notas con color, título y bordes redondeados:
\begin{lstlisting}
\begin{tcolorbox}[colback=blue!5, colframe=blue!60!black, title=Definicion]
  Una funcion es par si f(-x)=f(x).
\end{tcolorbox}
\end{lstlisting}
\produce
\begin{tcolorbox}[colback=blue!5,colframe=blue!60!black,title=Definición,
                  fonttitle=\bfseries]
  Una función es \textbf{par} si $f(-x)=f(x)$ para todo $x$ del dominio.
\end{tcolorbox}
\noindent \pkg{tcolorbox} es muy potente: cajas que se parten entre páginas,
\enquote{teoremas} con estilo, cajas con dos columnas, etc. La alternativa más
liviana es \pkg{mdframed}.

% =====================================================================
\section{Buenas prácticas}
% =====================================================================
\begin{itemize}[noitemsep,leftmargin=*]
  \item Tablas con \pkg{booktabs}; \textbf{sin} líneas verticales ni
        \C{hline}\C{hline}.
  \item \pkg{tabularx} para ancho total fijo; columna \texttt{S} para alinear
        decimales.
  \item Tabla con epígrafe $\to$ entorno \texttt{table} + \C{caption} + \C{label};
        citá con \C{cref}.
  \item Listas: ajustá con \pkg{enumitem}, no con \C{vspace} a mano.
  \item Cajas de color para \emph{resaltar con sentido}, sin abusar; \pkg{tcolorbox}
        para las elaboradas.
\end{itemize}

% =====================================================================
\section*{Ver también}
% =====================================================================
\addcontentsline{toc}{section}{Ver también}
\begin{itemize}[noitemsep,leftmargin=*]
  \item \emph{Estándares} — el canon de tablas (booktabs) y anti-patrones.
  \item \emph{Matemática} — entornos de teorema (otra forma de \enquote{caja}).
  \item Manuales: \texttt{texdoc booktabs}, \texttt{texdoc enumitem},
        \texttt{texdoc tcolorbox}.
\end{itemize}

\end{document}
